% Generated from the Markdown source -- edit the .md, not this file.
% Source of truth: projects/Project-Handbook.md
\documentclass[11pt]{article}

% Shared notes settings for Applied Machine Learning lectures (UPES).
% Keep this file free of \documentclass to allow reuse via \input.

\usepackage[utf8]{inputenc}
\usepackage[T1]{fontenc}
\usepackage{lmodern}          % scalable T1 fonts; without this pdflatex falls back
\input{glyphtounicode}        % to bitmapped EC Type 3 and the PDFs are neither
\pdfgentounicode=1            % searchable nor screen-reader friendly
\usepackage[a4paper,margin=1in]{geometry}
\usepackage{hyperref}
\usepackage{xurl}
\usepackage{booktabs}
\usepackage{amsmath}
\usepackage{amssymb}
\usepackage{listings}
\usepackage{xcolor}
\usepackage{enumitem}
\usepackage{parskip}

% Code listings (Python-oriented for this subject).
\lstset{
  language=Python,
  basicstyle=\ttfamily\small,
  keywordstyle=\color{blue},
  commentstyle=\color{gray},
  stringstyle=\color{teal},
  showstringspaces=false,
  columns=fullflexible,
  frame=single,
  framerule=0.2pt,
  breaklines=true
}

\setlist[itemize]{topsep=0.3em,itemsep=0.2em}
\setlist[enumerate]{topsep=0.3em,itemsep=0.2em}

% Simple per-section source line.
\newcommand{\notesources}[1]{\par\noindent{\small\color{gray}\textit{Sources:} \sloppy #1\par}}

% Unit-wise source strings for this subject (used by slides + notes).
% Path is relative to the lecture `latex/` folder (the compilation CWD).
\IfFileExists{../../../_shared/aml-sources.tex}{%
  % Source strings for Applied Machine Learning (CSAI2017P) — used by slides + notes.
% Prescribed texts and reference books are taken from the course syllabus.
% Support texts are the author-hosted open-access editions:
%   statlearning.com            (ISL with Python)
%   probml.github.io/pml-book   (Murphy, Probabilistic Machine Learning)
%   d2l.ai                      (Dive into Deep Learning)
%   mml-book.github.io          (Mathematics for Machine Learning)

\newcommand{\AMLRepoURL}{https://github.com/tali7c/Applied-Machine-Learning}

% --- prescribed -------------------------------------------------------------
\newcommand{\AMLPrescribed}{%
M\"uller and Guido, \textit{Introduction to Machine Learning with Python}, O'Reilly, 2016.;%
Bishop, \textit{Pattern Recognition and Machine Learning}, Springer, 2006.;%
Murphy, \textit{Machine Learning: A Probabilistic Perspective}, MIT Press, 2012.;%
Han, Kamber and Pei, \textit{Data Mining: Concepts and Techniques} (3e), Morgan Kaufmann, 2012.%
}

% --- unit-wise ---------------------------------------------------------------
\newcommand{\AMLUnitISources}{%
M\"uller and Guido, \textit{Introduction to Machine Learning with Python}, Ch. 1.;%
Bishop, \textit{Pattern Recognition and Machine Learning}, Ch. 1.;%
Murphy, \textit{Probabilistic Machine Learning: An Introduction}, MIT Press, 2022, Ch. 1.;%
Mitchell, \textit{Machine Learning}, McGraw Hill, 1997, Ch. 1.;%
Scikit-learn documentation, accessed 2026-08-04.%
}

\newcommand{\AMLUnitIISources}{%
Bishop, \textit{Pattern Recognition and Machine Learning}, \S1.5, \S4.3, \S7.1.;%
Zhang, Lipton, Li and Smola, \textit{Dive into Deep Learning}, \S3.4.;%
Shalev-Shwartz and Ben-David, \textit{Understanding Machine Learning}, Ch. 15.;%
Murphy, \textit{Probabilistic Machine Learning: An Introduction}, Ch. 5.%
}

\newcommand{\AMLUnitIIISources}{%
Zhang, Lipton, Li and Smola, \textit{Dive into Deep Learning}, Ch. 12 (Optimization Algorithms).;%
Deisenroth, Faisal and Ong, \textit{Mathematics for Machine Learning}, Ch. 7.;%
Kingma and Ba, ``Adam: A Method for Stochastic Optimization,'' ICLR 2015.;%
Ruder, ``An overview of gradient descent optimization algorithms,'' arXiv:1609.04747, 2016.%
}

\newcommand{\AMLUnitIVSources}{%
Han, Kamber and Pei, \textit{Data Mining: Concepts and Techniques} (3e), Ch. 3.;%
M\"uller and Guido, \textit{Introduction to Machine Learning with Python}, Ch. 3--4.;%
James, Witten, Hastie, Tibshirani and Taylor, \textit{An Introduction to Statistical Learning with Python}, Ch. 5--6, 12.;%
Scikit-learn documentation (preprocessing, model selection), accessed 2026-08-04.%
}

\newcommand{\AMLUnitVSources}{%
James et al., \textit{An Introduction to Statistical Learning with Python}, Ch. 3--6.;%
Bishop, \textit{Pattern Recognition and Machine Learning}, Ch. 3.;%
Deisenroth, Faisal and Ong, \textit{Mathematics for Machine Learning}, Ch. 9.;%
M\"uller and Guido, \textit{Introduction to Machine Learning with Python}, \S2.3.%
}

\newcommand{\AMLUnitVISources}{%
Han, Kamber and Pei, \textit{Data Mining: Concepts and Techniques} (3e), Ch. 8--9.;%
James et al., \textit{An Introduction to Statistical Learning with Python}, Ch. 8--9.;%
Bishop, \textit{Pattern Recognition and Machine Learning}, Ch. 4--5, 7, 14.;%
Zhang et al., \textit{Dive into Deep Learning}, Ch. 5.%
}

\newcommand{\AMLUnitVIISources}{%
Han, Kamber and Pei, \textit{Data Mining: Concepts and Techniques} (3e), Ch. 10.;%
James et al., \textit{An Introduction to Statistical Learning with Python}, Ch. 12.;%
M\"uller and Guido, \textit{Introduction to Machine Learning with Python}, \S3.5.;%
Ester, Kriegel, Sander and Xu, ``A density-based algorithm for discovering clusters (DBSCAN),'' KDD 1996.%
}

\newcommand{\AMLLabSources}{%
M\"uller and Guido, \textit{Introduction to Machine Learning with Python}, companion notebooks.;%
McKinney, \textit{Python for Data Analysis} (3e), O'Reilly, 2022.;%
Scikit-learn user guide, accessed 2026-08-04.;%
Matplotlib and Seaborn documentation, accessed 2026-08-04.%
}
%
}{%
  % Faculty copies live under `private/` so their compile CWD is different.
  \IfFileExists{../../../../public/_shared/aml-sources.tex}{%
    % Source strings for Applied Machine Learning (CSAI2017P) — used by slides + notes.
% Prescribed texts and reference books are taken from the course syllabus.
% Support texts are the author-hosted open-access editions:
%   statlearning.com            (ISL with Python)
%   probml.github.io/pml-book   (Murphy, Probabilistic Machine Learning)
%   d2l.ai                      (Dive into Deep Learning)
%   mml-book.github.io          (Mathematics for Machine Learning)

\newcommand{\AMLRepoURL}{https://github.com/tali7c/Applied-Machine-Learning}

% --- prescribed -------------------------------------------------------------
\newcommand{\AMLPrescribed}{%
M\"uller and Guido, \textit{Introduction to Machine Learning with Python}, O'Reilly, 2016.;%
Bishop, \textit{Pattern Recognition and Machine Learning}, Springer, 2006.;%
Murphy, \textit{Machine Learning: A Probabilistic Perspective}, MIT Press, 2012.;%
Han, Kamber and Pei, \textit{Data Mining: Concepts and Techniques} (3e), Morgan Kaufmann, 2012.%
}

% --- unit-wise ---------------------------------------------------------------
\newcommand{\AMLUnitISources}{%
M\"uller and Guido, \textit{Introduction to Machine Learning with Python}, Ch. 1.;%
Bishop, \textit{Pattern Recognition and Machine Learning}, Ch. 1.;%
Murphy, \textit{Probabilistic Machine Learning: An Introduction}, MIT Press, 2022, Ch. 1.;%
Mitchell, \textit{Machine Learning}, McGraw Hill, 1997, Ch. 1.;%
Scikit-learn documentation, accessed 2026-08-04.%
}

\newcommand{\AMLUnitIISources}{%
Bishop, \textit{Pattern Recognition and Machine Learning}, \S1.5, \S4.3, \S7.1.;%
Zhang, Lipton, Li and Smola, \textit{Dive into Deep Learning}, \S3.4.;%
Shalev-Shwartz and Ben-David, \textit{Understanding Machine Learning}, Ch. 15.;%
Murphy, \textit{Probabilistic Machine Learning: An Introduction}, Ch. 5.%
}

\newcommand{\AMLUnitIIISources}{%
Zhang, Lipton, Li and Smola, \textit{Dive into Deep Learning}, Ch. 12 (Optimization Algorithms).;%
Deisenroth, Faisal and Ong, \textit{Mathematics for Machine Learning}, Ch. 7.;%
Kingma and Ba, ``Adam: A Method for Stochastic Optimization,'' ICLR 2015.;%
Ruder, ``An overview of gradient descent optimization algorithms,'' arXiv:1609.04747, 2016.%
}

\newcommand{\AMLUnitIVSources}{%
Han, Kamber and Pei, \textit{Data Mining: Concepts and Techniques} (3e), Ch. 3.;%
M\"uller and Guido, \textit{Introduction to Machine Learning with Python}, Ch. 3--4.;%
James, Witten, Hastie, Tibshirani and Taylor, \textit{An Introduction to Statistical Learning with Python}, Ch. 5--6, 12.;%
Scikit-learn documentation (preprocessing, model selection), accessed 2026-08-04.%
}

\newcommand{\AMLUnitVSources}{%
James et al., \textit{An Introduction to Statistical Learning with Python}, Ch. 3--6.;%
Bishop, \textit{Pattern Recognition and Machine Learning}, Ch. 3.;%
Deisenroth, Faisal and Ong, \textit{Mathematics for Machine Learning}, Ch. 9.;%
M\"uller and Guido, \textit{Introduction to Machine Learning with Python}, \S2.3.%
}

\newcommand{\AMLUnitVISources}{%
Han, Kamber and Pei, \textit{Data Mining: Concepts and Techniques} (3e), Ch. 8--9.;%
James et al., \textit{An Introduction to Statistical Learning with Python}, Ch. 8--9.;%
Bishop, \textit{Pattern Recognition and Machine Learning}, Ch. 4--5, 7, 14.;%
Zhang et al., \textit{Dive into Deep Learning}, Ch. 5.%
}

\newcommand{\AMLUnitVIISources}{%
Han, Kamber and Pei, \textit{Data Mining: Concepts and Techniques} (3e), Ch. 10.;%
James et al., \textit{An Introduction to Statistical Learning with Python}, Ch. 12.;%
M\"uller and Guido, \textit{Introduction to Machine Learning with Python}, \S3.5.;%
Ester, Kriegel, Sander and Xu, ``A density-based algorithm for discovering clusters (DBSCAN),'' KDD 1996.%
}

\newcommand{\AMLLabSources}{%
M\"uller and Guido, \textit{Introduction to Machine Learning with Python}, companion notebooks.;%
McKinney, \textit{Python for Data Analysis} (3e), O'Reilly, 2022.;%
Scikit-learn user guide, accessed 2026-08-04.;%
Matplotlib and Seaborn documentation, accessed 2026-08-04.%
}
%
  }{%
    % Question papers live at `private/<family>/<instrument>/`.
    \IfFileExists{../../../public/_shared/aml-sources.tex}{%
      % Source strings for Applied Machine Learning (CSAI2017P) — used by slides + notes.
% Prescribed texts and reference books are taken from the course syllabus.
% Support texts are the author-hosted open-access editions:
%   statlearning.com            (ISL with Python)
%   probml.github.io/pml-book   (Murphy, Probabilistic Machine Learning)
%   d2l.ai                      (Dive into Deep Learning)
%   mml-book.github.io          (Mathematics for Machine Learning)

\newcommand{\AMLRepoURL}{https://github.com/tali7c/Applied-Machine-Learning}

% --- prescribed -------------------------------------------------------------
\newcommand{\AMLPrescribed}{%
M\"uller and Guido, \textit{Introduction to Machine Learning with Python}, O'Reilly, 2016.;%
Bishop, \textit{Pattern Recognition and Machine Learning}, Springer, 2006.;%
Murphy, \textit{Machine Learning: A Probabilistic Perspective}, MIT Press, 2012.;%
Han, Kamber and Pei, \textit{Data Mining: Concepts and Techniques} (3e), Morgan Kaufmann, 2012.%
}

% --- unit-wise ---------------------------------------------------------------
\newcommand{\AMLUnitISources}{%
M\"uller and Guido, \textit{Introduction to Machine Learning with Python}, Ch. 1.;%
Bishop, \textit{Pattern Recognition and Machine Learning}, Ch. 1.;%
Murphy, \textit{Probabilistic Machine Learning: An Introduction}, MIT Press, 2022, Ch. 1.;%
Mitchell, \textit{Machine Learning}, McGraw Hill, 1997, Ch. 1.;%
Scikit-learn documentation, accessed 2026-08-04.%
}

\newcommand{\AMLUnitIISources}{%
Bishop, \textit{Pattern Recognition and Machine Learning}, \S1.5, \S4.3, \S7.1.;%
Zhang, Lipton, Li and Smola, \textit{Dive into Deep Learning}, \S3.4.;%
Shalev-Shwartz and Ben-David, \textit{Understanding Machine Learning}, Ch. 15.;%
Murphy, \textit{Probabilistic Machine Learning: An Introduction}, Ch. 5.%
}

\newcommand{\AMLUnitIIISources}{%
Zhang, Lipton, Li and Smola, \textit{Dive into Deep Learning}, Ch. 12 (Optimization Algorithms).;%
Deisenroth, Faisal and Ong, \textit{Mathematics for Machine Learning}, Ch. 7.;%
Kingma and Ba, ``Adam: A Method for Stochastic Optimization,'' ICLR 2015.;%
Ruder, ``An overview of gradient descent optimization algorithms,'' arXiv:1609.04747, 2016.%
}

\newcommand{\AMLUnitIVSources}{%
Han, Kamber and Pei, \textit{Data Mining: Concepts and Techniques} (3e), Ch. 3.;%
M\"uller and Guido, \textit{Introduction to Machine Learning with Python}, Ch. 3--4.;%
James, Witten, Hastie, Tibshirani and Taylor, \textit{An Introduction to Statistical Learning with Python}, Ch. 5--6, 12.;%
Scikit-learn documentation (preprocessing, model selection), accessed 2026-08-04.%
}

\newcommand{\AMLUnitVSources}{%
James et al., \textit{An Introduction to Statistical Learning with Python}, Ch. 3--6.;%
Bishop, \textit{Pattern Recognition and Machine Learning}, Ch. 3.;%
Deisenroth, Faisal and Ong, \textit{Mathematics for Machine Learning}, Ch. 9.;%
M\"uller and Guido, \textit{Introduction to Machine Learning with Python}, \S2.3.%
}

\newcommand{\AMLUnitVISources}{%
Han, Kamber and Pei, \textit{Data Mining: Concepts and Techniques} (3e), Ch. 8--9.;%
James et al., \textit{An Introduction to Statistical Learning with Python}, Ch. 8--9.;%
Bishop, \textit{Pattern Recognition and Machine Learning}, Ch. 4--5, 7, 14.;%
Zhang et al., \textit{Dive into Deep Learning}, Ch. 5.%
}

\newcommand{\AMLUnitVIISources}{%
Han, Kamber and Pei, \textit{Data Mining: Concepts and Techniques} (3e), Ch. 10.;%
James et al., \textit{An Introduction to Statistical Learning with Python}, Ch. 12.;%
M\"uller and Guido, \textit{Introduction to Machine Learning with Python}, \S3.5.;%
Ester, Kriegel, Sander and Xu, ``A density-based algorithm for discovering clusters (DBSCAN),'' KDD 1996.%
}

\newcommand{\AMLLabSources}{%
M\"uller and Guido, \textit{Introduction to Machine Learning with Python}, companion notebooks.;%
McKinney, \textit{Python for Data Analysis} (3e), O'Reilly, 2022.;%
Scikit-learn user guide, accessed 2026-08-04.;%
Matplotlib and Seaborn documentation, accessed 2026-08-04.%
}
%
    }{%
      % Marking schemes live at `private/solutions/<family>/<id>/latex/`.
      \IfFileExists{../../../../../public/_shared/aml-sources.tex}{%
        % Source strings for Applied Machine Learning (CSAI2017P) — used by slides + notes.
% Prescribed texts and reference books are taken from the course syllabus.
% Support texts are the author-hosted open-access editions:
%   statlearning.com            (ISL with Python)
%   probml.github.io/pml-book   (Murphy, Probabilistic Machine Learning)
%   d2l.ai                      (Dive into Deep Learning)
%   mml-book.github.io          (Mathematics for Machine Learning)

\newcommand{\AMLRepoURL}{https://github.com/tali7c/Applied-Machine-Learning}

% --- prescribed -------------------------------------------------------------
\newcommand{\AMLPrescribed}{%
M\"uller and Guido, \textit{Introduction to Machine Learning with Python}, O'Reilly, 2016.;%
Bishop, \textit{Pattern Recognition and Machine Learning}, Springer, 2006.;%
Murphy, \textit{Machine Learning: A Probabilistic Perspective}, MIT Press, 2012.;%
Han, Kamber and Pei, \textit{Data Mining: Concepts and Techniques} (3e), Morgan Kaufmann, 2012.%
}

% --- unit-wise ---------------------------------------------------------------
\newcommand{\AMLUnitISources}{%
M\"uller and Guido, \textit{Introduction to Machine Learning with Python}, Ch. 1.;%
Bishop, \textit{Pattern Recognition and Machine Learning}, Ch. 1.;%
Murphy, \textit{Probabilistic Machine Learning: An Introduction}, MIT Press, 2022, Ch. 1.;%
Mitchell, \textit{Machine Learning}, McGraw Hill, 1997, Ch. 1.;%
Scikit-learn documentation, accessed 2026-08-04.%
}

\newcommand{\AMLUnitIISources}{%
Bishop, \textit{Pattern Recognition and Machine Learning}, \S1.5, \S4.3, \S7.1.;%
Zhang, Lipton, Li and Smola, \textit{Dive into Deep Learning}, \S3.4.;%
Shalev-Shwartz and Ben-David, \textit{Understanding Machine Learning}, Ch. 15.;%
Murphy, \textit{Probabilistic Machine Learning: An Introduction}, Ch. 5.%
}

\newcommand{\AMLUnitIIISources}{%
Zhang, Lipton, Li and Smola, \textit{Dive into Deep Learning}, Ch. 12 (Optimization Algorithms).;%
Deisenroth, Faisal and Ong, \textit{Mathematics for Machine Learning}, Ch. 7.;%
Kingma and Ba, ``Adam: A Method for Stochastic Optimization,'' ICLR 2015.;%
Ruder, ``An overview of gradient descent optimization algorithms,'' arXiv:1609.04747, 2016.%
}

\newcommand{\AMLUnitIVSources}{%
Han, Kamber and Pei, \textit{Data Mining: Concepts and Techniques} (3e), Ch. 3.;%
M\"uller and Guido, \textit{Introduction to Machine Learning with Python}, Ch. 3--4.;%
James, Witten, Hastie, Tibshirani and Taylor, \textit{An Introduction to Statistical Learning with Python}, Ch. 5--6, 12.;%
Scikit-learn documentation (preprocessing, model selection), accessed 2026-08-04.%
}

\newcommand{\AMLUnitVSources}{%
James et al., \textit{An Introduction to Statistical Learning with Python}, Ch. 3--6.;%
Bishop, \textit{Pattern Recognition and Machine Learning}, Ch. 3.;%
Deisenroth, Faisal and Ong, \textit{Mathematics for Machine Learning}, Ch. 9.;%
M\"uller and Guido, \textit{Introduction to Machine Learning with Python}, \S2.3.%
}

\newcommand{\AMLUnitVISources}{%
Han, Kamber and Pei, \textit{Data Mining: Concepts and Techniques} (3e), Ch. 8--9.;%
James et al., \textit{An Introduction to Statistical Learning with Python}, Ch. 8--9.;%
Bishop, \textit{Pattern Recognition and Machine Learning}, Ch. 4--5, 7, 14.;%
Zhang et al., \textit{Dive into Deep Learning}, Ch. 5.%
}

\newcommand{\AMLUnitVIISources}{%
Han, Kamber and Pei, \textit{Data Mining: Concepts and Techniques} (3e), Ch. 10.;%
James et al., \textit{An Introduction to Statistical Learning with Python}, Ch. 12.;%
M\"uller and Guido, \textit{Introduction to Machine Learning with Python}, \S3.5.;%
Ester, Kriegel, Sander and Xu, ``A density-based algorithm for discovering clusters (DBSCAN),'' KDD 1996.%
}

\newcommand{\AMLLabSources}{%
M\"uller and Guido, \textit{Introduction to Machine Learning with Python}, companion notebooks.;%
McKinney, \textit{Python for Data Analysis} (3e), O'Reilly, 2022.;%
Scikit-learn user guide, accessed 2026-08-04.;%
Matplotlib and Seaborn documentation, accessed 2026-08-04.%
}
%
      }{%
        % Source strings for Applied Machine Learning (CSAI2017P) — used by slides + notes.
% Prescribed texts and reference books are taken from the course syllabus.
% Support texts are the author-hosted open-access editions:
%   statlearning.com            (ISL with Python)
%   probml.github.io/pml-book   (Murphy, Probabilistic Machine Learning)
%   d2l.ai                      (Dive into Deep Learning)
%   mml-book.github.io          (Mathematics for Machine Learning)

\newcommand{\AMLRepoURL}{https://github.com/tali7c/Applied-Machine-Learning}

% --- prescribed -------------------------------------------------------------
\newcommand{\AMLPrescribed}{%
M\"uller and Guido, \textit{Introduction to Machine Learning with Python}, O'Reilly, 2016.;%
Bishop, \textit{Pattern Recognition and Machine Learning}, Springer, 2006.;%
Murphy, \textit{Machine Learning: A Probabilistic Perspective}, MIT Press, 2012.;%
Han, Kamber and Pei, \textit{Data Mining: Concepts and Techniques} (3e), Morgan Kaufmann, 2012.%
}

% --- unit-wise ---------------------------------------------------------------
\newcommand{\AMLUnitISources}{%
M\"uller and Guido, \textit{Introduction to Machine Learning with Python}, Ch. 1.;%
Bishop, \textit{Pattern Recognition and Machine Learning}, Ch. 1.;%
Murphy, \textit{Probabilistic Machine Learning: An Introduction}, MIT Press, 2022, Ch. 1.;%
Mitchell, \textit{Machine Learning}, McGraw Hill, 1997, Ch. 1.;%
Scikit-learn documentation, accessed 2026-08-04.%
}

\newcommand{\AMLUnitIISources}{%
Bishop, \textit{Pattern Recognition and Machine Learning}, \S1.5, \S4.3, \S7.1.;%
Zhang, Lipton, Li and Smola, \textit{Dive into Deep Learning}, \S3.4.;%
Shalev-Shwartz and Ben-David, \textit{Understanding Machine Learning}, Ch. 15.;%
Murphy, \textit{Probabilistic Machine Learning: An Introduction}, Ch. 5.%
}

\newcommand{\AMLUnitIIISources}{%
Zhang, Lipton, Li and Smola, \textit{Dive into Deep Learning}, Ch. 12 (Optimization Algorithms).;%
Deisenroth, Faisal and Ong, \textit{Mathematics for Machine Learning}, Ch. 7.;%
Kingma and Ba, ``Adam: A Method for Stochastic Optimization,'' ICLR 2015.;%
Ruder, ``An overview of gradient descent optimization algorithms,'' arXiv:1609.04747, 2016.%
}

\newcommand{\AMLUnitIVSources}{%
Han, Kamber and Pei, \textit{Data Mining: Concepts and Techniques} (3e), Ch. 3.;%
M\"uller and Guido, \textit{Introduction to Machine Learning with Python}, Ch. 3--4.;%
James, Witten, Hastie, Tibshirani and Taylor, \textit{An Introduction to Statistical Learning with Python}, Ch. 5--6, 12.;%
Scikit-learn documentation (preprocessing, model selection), accessed 2026-08-04.%
}

\newcommand{\AMLUnitVSources}{%
James et al., \textit{An Introduction to Statistical Learning with Python}, Ch. 3--6.;%
Bishop, \textit{Pattern Recognition and Machine Learning}, Ch. 3.;%
Deisenroth, Faisal and Ong, \textit{Mathematics for Machine Learning}, Ch. 9.;%
M\"uller and Guido, \textit{Introduction to Machine Learning with Python}, \S2.3.%
}

\newcommand{\AMLUnitVISources}{%
Han, Kamber and Pei, \textit{Data Mining: Concepts and Techniques} (3e), Ch. 8--9.;%
James et al., \textit{An Introduction to Statistical Learning with Python}, Ch. 8--9.;%
Bishop, \textit{Pattern Recognition and Machine Learning}, Ch. 4--5, 7, 14.;%
Zhang et al., \textit{Dive into Deep Learning}, Ch. 5.%
}

\newcommand{\AMLUnitVIISources}{%
Han, Kamber and Pei, \textit{Data Mining: Concepts and Techniques} (3e), Ch. 10.;%
James et al., \textit{An Introduction to Statistical Learning with Python}, Ch. 12.;%
M\"uller and Guido, \textit{Introduction to Machine Learning with Python}, \S3.5.;%
Ester, Kriegel, Sander and Xu, ``A density-based algorithm for discovering clusters (DBSCAN),'' KDD 1996.%
}

\newcommand{\AMLLabSources}{%
M\"uller and Guido, \textit{Introduction to Machine Learning with Python}, companion notebooks.;%
McKinney, \textit{Python for Data Analysis} (3e), O'Reilly, 2022.;%
Scikit-learn user guide, accessed 2026-08-04.;%
Matplotlib and Seaborn documentation, accessed 2026-08-04.%
}
%
      }%
    }%
  }%
}


\usepackage{longtable}
\usepackage{array}

\DeclareUnicodeCharacter{2713}{\ensuremath{\checkmark}}
\DeclareUnicodeCharacter{2192}{\ensuremath{\rightarrow}}
\DeclareUnicodeCharacter{2212}{\ensuremath{-}}
\DeclareUnicodeCharacter{2265}{\ensuremath{\geq}}
\DeclareUnicodeCharacter{2264}{\ensuremath{\leq}}
\DeclareUnicodeCharacter{2248}{\ensuremath{\approx}}
\DeclareUnicodeCharacter{00D7}{\ensuremath{\times}}
\DeclareUnicodeCharacter{00B1}{\ensuremath{\pm}}
\DeclareUnicodeCharacter{00B2}{\ensuremath{^2}}
\DeclareUnicodeCharacter{00B3}{\ensuremath{^3}}
\DeclareUnicodeCharacter{2500}{\textbar}
\DeclareUnicodeCharacter{2502}{\textbar}
\DeclareUnicodeCharacter{251C}{\textbar}
\DeclareUnicodeCharacter{2514}{\textbar}

% Section numbers come from the Markdown headings themselves (the documents
% cross-reference them as "handbook 3.3"), so LaTeX must not add its own.
\setcounter{secnumdepth}{0}
\setcounter{tocdepth}{2}

% pandoc fragments rely on these being defined by its default template.
\providecommand{\tightlist}{\setlength{\itemsep}{0pt}\setlength{\parskip}{0pt}}
\providecommand{\passthrough}[1]{#1}

% Dataset URLs are long and unbreakable; xurl (from notes-common) plus a
% tolerant line-breaker keeps them inside the text block.
\sloppy
% Coloured rather than boxed links: legible on screen, clean in print.
\hypersetup{colorlinks=true, linkcolor=UPESBlueD, urlcolor=UPESBlueD,
            citecolor=UPESBlueD, linktoc=all}

\title{Project Handbook\\[4pt]\large Rules that apply to every theme}
\author{Dr.\ Tofik Ali \\ School of Computer Science, UPES Dehradun}
\date{Applied Machine Learning (CSAI2017P) \quad\textbullet\quad B.Tech CSE (AI \& ML), Semester III \\[2pt] Autumn 2026, AY 2026-27}

\begin{document}
\maketitle
\tableofcontents
\newpage

\hypertarget{aml-project-handbook-rules-that-apply-to-every-theme}{%
\section{AML Project Handbook --- rules that apply to every theme}\label{aml-project-handbook-rules-that-apply-to-every-theme}}

Applied Machine Learning (CSAI2017P) · Autumn 2026 · Dr.~Tofik Ali

Read this once. It applies to all eight themes P01--P08. The theme folder tells you
\emph{what} to build; this handbook tells you \emph{how it is submitted and marked}.



\hypertarget{the-component-in-one-table}{%
\subsection{1. The component in one table}\label{the-component-in-one-table}}

{\footnotesize
\begin{longtable}[]{@{}llll@{}}
\toprule
\begin{minipage}[b]{0.22\columnwidth}\raggedright
Sub-component\strut
\end{minipage} & \begin{minipage}[b]{0.22\columnwidth}\raggedright
Marks\strut
\end{minipage} & \begin{minipage}[b]{0.22\columnwidth}\raggedright
Granularity\strut
\end{minipage} & \begin{minipage}[b]{0.22\columnwidth}\raggedright
Basis\strut
\end{minipage}\tabularnewline
\midrule
\endhead
\begin{minipage}[t]{0.22\columnwidth}\raggedright
\textbf{Project quiz} --- a written paper, one set per theme\strut
\end{minipage} & \begin{minipage}[t]{0.22\columnwidth}\raggedright
\textbf{10}\strut
\end{minipage} & \begin{minipage}[t]{0.22\columnwidth}\raggedright
\textbf{per student}\strut
\end{minipage} & \begin{minipage}[t]{0.22\columnwidth}\raggedright
your own understanding of your group's work\strut
\end{minipage}\tabularnewline
\begin{minipage}[t]{0.22\columnwidth}\raggedright
\textbf{Repository submission} --- the state and the history of the repo\strut
\end{minipage} & \begin{minipage}[t]{0.22\columnwidth}\raggedright
\textbf{5}\strut
\end{minipage} & \begin{minipage}[t]{0.22\columnwidth}\raggedright
\textbf{per group}\strut
\end{minipage} & \begin{minipage}[t]{0.22\columnwidth}\raggedright
identical mark for every member\strut
\end{minipage}\tabularnewline
\bottomrule
\end{longtable}
}

Total 15 of the 100-mark course scale.

\textbf{The 5 group marks are identical for every member of the group. They are not
split by contribution.} All individual differentiation happens in the 10-mark
written quiz, which every member sits separately. A member who did nothing still
receives the group's 5 marks --- and then faces a paper about work they did not do.
That is the design. Choose your group knowing it.

\textbf{Both parts are compulsory.} A student who does not sit the quiz scores 0 of 10
for it; a group that does not submit a repository link scores 0 of 5.



\hypertarget{groups}{%
\subsection{2. Groups}\label{groups}}

\begin{itemize}
\tightlist
\item
  \textbf{4 to 8 students}, self-formed, from the same section.
\item
  \textbf{At most two groups per theme.} Sign-up is first come, first served on the
  sheet; once a theme has two groups it is closed.
\item
  One member is the \textbf{repo owner} --- they create the repository and add everyone
  else as a collaborator. This is an administrative role, not a leadership one,
  and it earns no extra marks.
\item
  Group membership is frozen at the end of Unit IV. Tell me before that if it
  changes; after that it does not change.
\end{itemize}



\hypertarget{the-github-repository-the-only-thing-you-submit}{%
\subsection{3. The GitHub repository --- the only thing you submit}\label{the-github-repository-the-only-thing-you-submit}}

At the end, you send me \textbf{one line: the repository URL.} Nothing else. No zip,
no email attachment, no printed report.

\hypertarget{creating-it}{%
\subsubsection{3.1 Creating it}\label{creating-it}}

\begin{itemize}
\tightlist
\item
  One \textbf{public} repository per group, on GitHub, named
  \texttt{aml-2026-\textless{}theme-code\textgreater{}-\textless{}group-number\textgreater{}} --- for example \texttt{aml-2026-p03-g2}.
\item
  The repo owner adds \textbf{every member as a collaborator}, and adds me as a
  collaborator too if you choose to keep it private (public is simpler; either
  is accepted).
\item
  Create it \textbf{by the end of Unit IV} --- at the start of the work, not at the end
  of the semester. The creation date is visible and it is part of what I look at.
\end{itemize}

\hypertarget{what-i-check-when-you-send-the-link}{%
\subsubsection{3.2 What I check when you send the link}\label{what-i-check-when-you-send-the-link}}

I open the repository and read \textbf{the commit history}, not just the final files.
Specifically:

{\footnotesize
\begin{longtable}[]{@{}ll@{}}
\toprule
\begin{minipage}[b]{0.47\columnwidth}\raggedright
I look at\strut
\end{minipage} & \begin{minipage}[b]{0.47\columnwidth}\raggedright
What I am checking\strut
\end{minipage}\tabularnewline
\midrule
\endhead
\begin{minipage}[t]{0.47\columnwidth}\raggedright
\texttt{git\ log} over time\strut
\end{minipage} & \begin{minipage}[t]{0.47\columnwidth}\raggedright
that the work was spread across the semester\strut
\end{minipage}\tabularnewline
\begin{minipage}[t]{0.47\columnwidth}\raggedright
who authored each commit\strut
\end{minipage} & \begin{minipage}[t]{0.47\columnwidth}\raggedright
that more than one person actually pushed\strut
\end{minipage}\tabularnewline
\begin{minipage}[t]{0.47\columnwidth}\raggedright
commit messages\strut
\end{minipage} & \begin{minipage}[t]{0.47\columnwidth}\raggedright
that they describe changes, not ``update'', ``final'', ``final2''\strut
\end{minipage}\tabularnewline
\begin{minipage}[t]{0.47\columnwidth}\raggedright
the two checkpoint tags\strut
\end{minipage} & \begin{minipage}[t]{0.47\columnwidth}\raggedright
that the checkpoints were real, not backdated\strut
\end{minipage}\tabularnewline
\begin{minipage}[t]{0.47\columnwidth}\raggedright
diffs at each checkpoint\strut
\end{minipage} & \begin{minipage}[t]{0.47\columnwidth}\raggedright
that the state then matches what you claimed then\strut
\end{minipage}\tabularnewline
\bottomrule
\end{longtable}
}

\hypertarget{the-commit-history-floor}{%
\subsubsection{3.3 The commit-history floor}\label{the-commit-history-floor}}

A repository that meets all of these is safe:

\begin{enumerate}
\def\labelenumi{\arabic{enumi}.}
\tightlist
\item
  \textbf{Created by the end of Unit IV} and first commit within that week.
\item
  \textbf{At least 20 commits} by the final submission.
\item
  Commits fall on \textbf{at least 6 distinct calendar weeks}.
\item
  \textbf{No more than 40 \% of the commits in the final week.}
\item
  \textbf{Every member has commits under their own GitHub account.} Use your own
  account, set \texttt{user.email} correctly, and do not let one person push everyone's
  work. This is a condition on the \emph{group's} repository, so a repository where
  one person pushed everything loses process marks for the whole group --- and the
  members who never committed will also find the quiz very hard.
\item
  Two annotated tags, \texttt{checkpoint-1} and \texttt{checkpoint-2}, pushed \textbf{at the time of
  each checkpoint}.
\end{enumerate}

\textbf{What fails.} A repository whose entire history is 1--3 commits in the last few
days before submission --- the whole project uploaded at once, files unorganised, one commit called
``project'' --- scores \textbf{at most 1.5 of 5} for the group component, however good the
code inside it is. This is stated in advance and applied without exception.

\textbf{Honest exceptions exist.} If a genuine problem broke your history --- a repo
recreated after an accident, a member's account issue --- tell me \emph{at the time},
not at submission. A note in the README written on the day it happened is
evidence; an explanation invented at the end is not.

\hypertarget{repository-layout}{%
\subsubsection{3.4 Repository layout}\label{repository-layout}}

Every group uses this structure. It is the same for all eight themes; only the
number and names of the notebooks vary, and your theme brief tells you which.

\begin{verbatim}
aml-2026-pNN-gM/
|-- README.md            # what the project is, how to run it, who is in the group
|-- requirements.txt     # pinned versions
|-- .gitignore
|-- data/
|   |-- raw/             # the downloaded file(s), unmodified  (see §4)
|   `-- processed/       # generated — usually gitignored
|-- notebooks/          # numbered in run order; your theme brief names them
|   |-- 01-eda.ipynb
|   |-- 02-preprocessing.ipynb
|   |-- 03-models.ipynb
|   `-- 04-evaluation.ipynb
|-- src/                 # reusable functions imported by the notebooks
|-- results/
|   |-- figures/
|   `-- metrics.csv      # every number you report, in one file
|-- report/
|   `-- report.pdf       # 4–6 pages
`-- REPRODUCIBILITY.md   # §3.6
\end{verbatim}

\hypertarget{the-readme-must-contain}{%
\subsubsection{3.5 The README must contain}\label{the-readme-must-contain}}

\begin{itemize}
\tightlist
\item
  The theme code and title, and the group number.
\item
  \textbf{Every member: name, enrolment number, GitHub username.}
\item
  One paragraph on the question the project answers.
\item
  The data source URL and the exact download steps.
\item
  How to run it: environment, then the commands, in order.
\item
  The headline result --- the metric, the number, and on which split.
\item
  A short \textbf{``who did what''} section. It does not change any mark. It exists so
  that the group agrees, in writing and in advance, on what each person will be
  able to answer for in the quiz.
\end{itemize}

\hypertarget{reproducibility.md-must-contain}{%
\subsubsection{3.6 REPRODUCIBILITY.md must contain}\label{reproducibility.md-must-contain}}

\begin{itemize}
\tightlist
\item
  Python version and how to recreate the environment.
\item
  Every random seed used, and where it is set.
\item
  The exact sequence of notebooks/scripts to run.
\item
  Expected runtime and hardware.
\item
  Any manual step (a browser download, a file placed by hand) spelled out.
\item
  Anything that does \textbf{not} reproduce exactly, and why.
\end{itemize}

I will attempt a clean-checkout run on at least a sample of repositories.
A project that cannot be re-run is a claim, not a result.



\hypertarget{data-on-the-upes-network}{%
\subsection{4. Data on the UPES network}\label{data-on-the-upes-network}}

The campus TLS appliance breaks Python downloads --- \texttt{fetch\_openml}, \texttt{ucimlrepo},
\texttt{yfinance}, \texttt{requests}, \texttt{pip} from some hosts. Details in
\texttt{../lab-assignments/Anchor-Datasets.pdf}.

So, for every theme:

\begin{itemize}
\tightlist
\item
  \textbf{Download the file once in a browser} (at home, on mobile data, or on campus ---
  browsers work), and commit it to \texttt{data/raw/} if it is under 50 MB.
\item
  If it is larger than 50 MB, do \textbf{not} commit it. Put it in \texttt{.gitignore},
  and write the download instructions in the README precisely enough that I can
  reproduce your \texttt{data/raw/} myself in under two minutes.
\item
  Mirrors of the raw files are placed on the LMS where licensing permits; check
  there first.
\item
  \textbf{Never disable certificate verification.} \texttt{verify=False} or an unverified SSL
  context anywhere in the repository costs marks. It is the wrong lesson.
\end{itemize}



\hypertarget{milestones}{%
\subsection{5. Milestones}\label{milestones}}

Anchored to units, not dates. Exact dates announced in class.

{\footnotesize
\begin{longtable}[]{@{}llll@{}}
\toprule
\begin{minipage}[b]{0.22\columnwidth}\raggedright
\#\strut
\end{minipage} & \begin{minipage}[b]{0.22\columnwidth}\raggedright
Milestone\strut
\end{minipage} & \begin{minipage}[b]{0.22\columnwidth}\raggedright
Anchored to\strut
\end{minipage} & \begin{minipage}[b]{0.22\columnwidth}\raggedright
Evidence\strut
\end{minipage}\tabularnewline
\midrule
\endhead
\begin{minipage}[t]{0.22\columnwidth}\raggedright
1\strut
\end{minipage} & \begin{minipage}[t]{0.22\columnwidth}\raggedright
Catalogue released\strut
\end{minipage} & \begin{minipage}[t]{0.22\columnwidth}\raggedright
during Unit IV\strut
\end{minipage} & \begin{minipage}[t]{0.22\columnwidth}\raggedright
---\strut
\end{minipage}\tabularnewline
\begin{minipage}[t]{0.22\columnwidth}\raggedright
2\strut
\end{minipage} & \begin{minipage}[t]{0.22\columnwidth}\raggedright
Groups formed, theme selected, \textbf{repo created}\strut
\end{minipage} & \begin{minipage}[t]{0.22\columnwidth}\raggedright
end of Unit IV\strut
\end{minipage} & \begin{minipage}[t]{0.22\columnwidth}\raggedright
repo URL on the sign-up sheet\strut
\end{minipage}\tabularnewline
\begin{minipage}[t]{0.22\columnwidth}\raggedright
3\strut
\end{minipage} & \begin{minipage}[t]{0.22\columnwidth}\raggedright
\textbf{Checkpoint 1} --- data loaded, one baseline number\strut
\end{minipage} & \begin{minipage}[t]{0.22\columnwidth}\raggedright
end of Unit V\strut
\end{minipage} & \begin{minipage}[t]{0.22\columnwidth}\raggedright
tag \texttt{checkpoint-1}\strut
\end{minipage}\tabularnewline
\begin{minipage}[t]{0.22\columnwidth}\raggedright
4\strut
\end{minipage} & \begin{minipage}[t]{0.22\columnwidth}\raggedright
\textbf{Checkpoint 2} --- full pipeline, initial results\strut
\end{minipage} & \begin{minipage}[t]{0.22\columnwidth}\raggedright
end of Unit VI\strut
\end{minipage} & \begin{minipage}[t]{0.22\columnwidth}\raggedright
tag \texttt{checkpoint-2}\strut
\end{minipage}\tabularnewline
\begin{minipage}[t]{0.22\columnwidth}\raggedright
5\strut
\end{minipage} & \begin{minipage}[t]{0.22\columnwidth}\raggedright
\textbf{Final submission} --- repo link sent\strut
\end{minipage} & \begin{minipage}[t]{0.22\columnwidth}\raggedright
end of Unit VII\strut
\end{minipage} & \begin{minipage}[t]{0.22\columnwidth}\raggedright
the URL\strut
\end{minipage}\tabularnewline
\begin{minipage}[t]{0.22\columnwidth}\raggedright
6\strut
\end{minipage} & \begin{minipage}[t]{0.22\columnwidth}\raggedright
\textbf{Project quiz} --- written, individual\strut
\end{minipage} & \begin{minipage}[t]{0.22\columnwidth}\raggedright
final lecture / lab-15 slot\strut
\end{minipage} & \begin{minipage}[t]{0.22\columnwidth}\raggedright
---\strut
\end{minipage}\tabularnewline
\bottomrule
\end{longtable}
}

\textbf{Checkpoint 1 is real, not a plan.} You show the data actually loaded --- shape,
head, a chart --- and one honest baseline number. Groups whose dataset turns out not
to exist are caught here, when there is still time to move. A checkpoint tag
pushed after the fact is worse than a missed checkpoint honestly declared.



\hypertarget{the-5-group-marks-rubric}{%
\subsection{6. The 5 group marks --- rubric}\label{the-5-group-marks-rubric}}

{\footnotesize
\begin{longtable}[]{@{}lll@{}}
\toprule
\begin{minipage}[b]{0.30\columnwidth}\raggedright
Criterion\strut
\end{minipage} & \begin{minipage}[b]{0.30\columnwidth}\raggedright
Marks\strut
\end{minipage} & \begin{minipage}[b]{0.30\columnwidth}\raggedright
Full marks looks like\strut
\end{minipage}\tabularnewline
\midrule
\endhead
\begin{minipage}[t]{0.30\columnwidth}\raggedright
\textbf{Process \& commit history}\strut
\end{minipage} & \begin{minipage}[t]{0.30\columnwidth}\raggedright
\textbf{1.0}\strut
\end{minipage} & \begin{minipage}[t]{0.30\columnwidth}\raggedright
meets all six floor conditions in §3.3 --- including commits from every member's own account; messages describe changes; both checkpoint tags pushed on time\strut
\end{minipage}\tabularnewline
\begin{minipage}[t]{0.30\columnwidth}\raggedright
\textbf{Completeness}\strut
\end{minipage} & \begin{minipage}[t]{0.30\columnwidth}\raggedright
\textbf{1.5}\strut
\end{minipage} & \begin{minipage}[t]{0.30\columnwidth}\raggedright
every item in the theme's ``Done'' list is present and demonstrated\strut
\end{minipage}\tabularnewline
\begin{minipage}[t]{0.30\columnwidth}\raggedright
\textbf{Method correctness}\strut
\end{minipage} & \begin{minipage}[t]{0.30\columnwidth}\raggedright
\textbf{1.0}\strut
\end{minipage} & \begin{minipage}[t]{0.30\columnwidth}\raggedright
no leakage; split appropriate to the data (stratified / chronological); metric fits the problem and the choice is argued\strut
\end{minipage}\tabularnewline
\begin{minipage}[t]{0.30\columnwidth}\raggedright
\textbf{Reproducibility}\strut
\end{minipage} & \begin{minipage}[t]{0.30\columnwidth}\raggedright
\textbf{1.0}\strut
\end{minipage} & \begin{minipage}[t]{0.30\columnwidth}\raggedright
runs from a clean checkout following REPRODUCIBILITY.md; seeds set; reported numbers reappear\strut
\end{minipage}\tabularnewline
\begin{minipage}[t]{0.30\columnwidth}\raggedright
\textbf{Documentation}\strut
\end{minipage} & \begin{minipage}[t]{0.30\columnwidth}\raggedright
\textbf{0.5}\strut
\end{minipage} & \begin{minipage}[t]{0.30\columnwidth}\raggedright
README and report are complete, honest about limitations, and readable\strut
\end{minipage}\tabularnewline
\begin{minipage}[t]{0.30\columnwidth}\raggedright
\strut
\end{minipage} & \begin{minipage}[t]{0.30\columnwidth}\raggedright
\textbf{5.0}\strut
\end{minipage} & \begin{minipage}[t]{0.30\columnwidth}\raggedright
\strut
\end{minipage}\tabularnewline
\bottomrule
\end{longtable}
}

The ``Further'' items in a theme brief are not needed for full marks. They are where
a strong group makes the quiz easy for itself.



\hypertarget{the-project-quiz-10-marks-written-individual}{%
\subsection{7. The project quiz --- 10 marks, written, individual}\label{the-project-quiz-10-marks-written-individual}}

\begin{itemize}
\item
  \textbf{A written paper. One set per theme}, so groups on different themes sit
  different papers. Both groups on the same theme sit the same set.
\item
  Conducted in the final lecture / lab slot. Closed-book unless announced
  otherwise; your own repository may be open on screen.
\item
  Every member sits it \textbf{separately}. Answers are your own.
\item
  The questions are \textbf{not pre-released}. What is pre-released is the \emph{kind} of
  question:

  \begin{enumerate}
  \def\labelenumi{\arabic{enumi}.}
  \tightlist
  \item
    Trace the path from the raw file to one number in your report.
  \item
    You reported metric M. Why M and not the obvious alternative?
  \item
    Here is a block of code from your repository. What does it do? What breaks
    if line \emph{k} is removed?
  \item
    Your result would change if X changed. Up or down, and why?
  \item
    What is the weakest part of your project, and what would you do with two
    more weeks?
  \end{enumerate}
\item
  \textbf{Marking is on understanding, not fluency.} ``I don't know, but I would find
  out by running Y and looking at Z'' earns more than a memorised paragraph that
  does not fit the question.
\item
  A contributor answers these comfortably. A passenger cannot, however polished
  the repository is.
\end{itemize}



\hypertarget{academic-integrity}{%
\subsection{8. Academic integrity}\label{academic-integrity}}

\begin{itemize}
\tightlist
\item
  Using AI tools, tutorials, Kaggle notebooks or Stack Overflow is \textbf{allowed and
  expected} --- provided you (a) say so in the README, citing what you used and
  where, and (b) can explain every line you kept. The quiz tests (b).
\item
  Copying another group's repository, or submitting a lab assignment as a project,
  is plagiarism and is reported.
\item
  The eight themes deliberately avoid every lab anchor dataset, so a re-submitted
  lab is immediately visible.
\end{itemize}



\hypertarget{getting-help}{%
\subsection{9. Getting help}\label{getting-help}}

Bring a specific question and your repository link. ``It doesn't work'' is not a
question I can answer; ``the chronological split gives an R² of −0.3 and I expected
it to be positive, here is the cell'' is.

Ask early. The gap between Checkpoint 1 and Checkpoint 2 is where projects are won.

\end{document}
